Als abgeschlossene beschränkte Teilmenge des $\R^2$ ist der Einheitskreis kompakt
\zusatzklammer {[[Kompaktheit/Satz von Heine-Borel/Fakt|Satz von Heine-Borel]]} {} {.} Die reelle Gerade $\R$ und $\R \setminus \{0\}$ sind hingegen nicht kompakt, da sie unbeschränkte Teilmengen von $\R$ sind. Also ist
\mathl{S^1 \not \cong \R,\, \R \setminus \{0\}}{.}

Die reelle Gerade ist zusammenhängend, wie aus dem
[[Reelle Analysis/Zwischenwertsatz/Fakt|Zwischenwertsatz]]
folgt. Dagegen ist $\R \setminus \{0\}$ nicht zusammenhängend, da man

\mavergleichskettedisp
{\vergleichskette
{\R_+
}
{ =} {(\R \setminus \{0\}) \cap [0, \infty]
}
{ =} {(\R \setminus \{0\}) \cap \, ]0, \infty]
}
{ } {
}
{ } {
}
}
{}{}{}
schreiben kann, sodass man eine sowohl offene als auch abgeschlossene nichtleere Teilmenge erhält. Also ist
\mathl{\R \,\not \cong \R \setminus \{0\}}{.}

 [[Kategorie:Latexseite]]
